\documentclass[amssymb,pra,12pt,aps,notitlepage]{revtex4-1}
\usepackage{mathptmx,amsmath, multirow, graphicx}
\usepackage{parskip}
\begin{document}

\title{Putnam POTD \today}
\maketitle

\section{Problem}
(1977 B5)
Prove that for $n\geq 2$,
\[
\overbrace{2^{2^{\cdots^{2}}}}^{\mbox{$n$ terms}} \equiv
\overbrace{2^{2^{\cdots^{2}}}}^{\mbox{$n-1$ terms}} \quad \pmod{n}.
\]

\section{Solution}
We redefine the variables for style. Let 
\[x_i = \overbrace{2^{2^{\cdots^{2}}}}^{\mbox{$i$ terms}}\]
for all $i \in \mathbb{N}$.

Our goal is to show that $\forall n \ge 2 \in \mathbb{Z}$, 
there exists some $m_n$ such that 
\[x_{m_n} = x_{m_n+1} = \dots \pmod{n}.\]

We proceed with strong induction on $n$. The base case 
is $n = 2$, in this case, 
$x_1 = 2, x_2 = 2^2 = 4, x_3 = 2^{2^2} = 16 ...$, 
so $m_1 = 1$, the statement holds for $n = 2$. 

The inductive hypothesis: assume this works for all integers $2, \dots, n-1$ 
for some integer $n \ge 3$, we want to show that there is also 
a desired integer $m_n$. 

We can split $n = 2^a b$ to divide and conquer using the Chinese Remainder Theorem,
considering the sequence mod $2^a$ and the sequence mod $b$, and deducing the 
behavior under mod $n$. 

Under mod $2^a$, for all $j \in \mathbb{N}$ such that $x_{j-1} \ge a$, $x_j = 0 \pmod{2^a}$,
this is because $x_j = 2^{x_{j-1}}$ and we need to ensure that 
$x_{j-1} \ge a$, so we set $m_a$ to be the smallest $j$ such $j$. 

Under mod $b$, we note that $\gcd(b,2) = 1$, so from Euler's theorem: 

\[2^{\phi(b)} = 1 \pmod{b}\]
then for $x_j = x_{j+1} \pmod{b}$, we only need to ensure that 
$2^{x_{j-1}} = 2^{x_{j}}\pmod{b}$, so we need 
\[x_{j-1} = x_j \pmod{\phi(b)}.\]

Here comes the punchline! Since $\phi(b) < b < n$, there exists 
$m_{\phi(b)}$ such that 
\[x_{m_{\phi(b)}} = x_{m_{\phi(b)+1}} = \dots \pmod{\phi(b)},\]
then 
we take $m_b = m_{\phi(b)+1}$, so that 
\[x_{m_b} = x_{m_b+1} = \dots \pmod{b
}.\]

Now define $m_n = \max(m_a,m_b)$, then 
\begin{align*}
    x_{m_n} &= x_{m_n+1} \pmod{2^a} \\ 
    x_{m_n} &= x_{m_n+1} \pmod{b} \\ 
\end{align*}
hence $2^a \mid x_{m_n+1} - x_{m_n}$ and $b \mid x_{m_n+1} - x_{m_n}$, and since $\gcd(2^a, b) = 1$, 
$n = 2^ab \mid x_{m_n+1} - x_{m_n}$
hence 
\[x_{m_n} = x_{m_n+1} = \dots \pmod{n}.\]


Furthermore, since $m_b$ is always larger than $m_a$, $m_n = m_b \le b < n$, so the original claim always holds.




\end{document}